\noindent\makebox[0pt][l]{\fontsize{13}{16}\selectfont 418}\hfill{\small\textbf{NACHLASS.}}\hfill\null\par\medskip

\begin{center}
Historischer Bericht\\
über die von Hofrat \textsc{Gauss} teils ausgeführten, teils geleiteten\\
Messungen im Königreich Hannover.
\end{center}

\par\noindent\hspace*{\parindent}Die verschiedenen Vermessungsarbeiten, über die ich hier einen kurzen historischen Bericht zu geben habe, sind zwar eng miteinander verknüpft, haben aber unterschiedliche Zwecke und einen unterschiedlichen Charakter. Am übersichtlichsten ist es, sie in der zeitlichen Reihenfolge der Aufträge zu ordnen, aufgrund deren ich sie übernommen habe.

\par\medskip
\begin{center}
I.\quad Die hannoversche Gradmessung.
\end{center}

\par\noindent\hspace*{\parindent}Durch einen Erlass vom 30. Juni 1820 wurde ich beauftragt, eine Gradmessung durch das Königreich Hannover auszuführen, als Erweiterung oder Fortsetzung einer ähnlichen Arbeit, die kurz zuvor in den dänischen Gebieten begonnen worden war. Unter dieser Bezeichnung versteht man die teils astronomischen, teils trigonometrischen Arbeiten, mit denen die Größe eines Meridiangrades in einem bekannten Längenmaß (Fuß, Toise usw.) bestimmt wird. Dazu wählt man zwei hinreichend weit voneinander entfernte Punkte auf demselben Meridian, bestimmt die Länge des zwischen ihnen liegenden Bogens in Fuß usw. durch ein zwischen ihnen angelegtes Dreiecksnetz und ermittelt durch astronomische Beobachtungen an den Endpunkten die Anzahl der Grade, Minuten und Sekunden, die diesem Bogen entsprechen. Daraus schließt man dann auf die Länge eines Grades.

\par\noindent\hspace*{\parindent}Daraus geht hervor, dass eine Gradmessung \textls[120]{als solche} nur den rein wissenschaftlichen Zweck hat, zur mathematischen Kenntnis der Verhältnisse des Erdsphäroids beizutragen. Diese erfordert nach den Ansprüchen des Jahrhunderts eine viel größere Genauigkeit und eine viel größere Zahl von Gradmessungen, als früher genügen mussten. Ferner geht daraus hervor, dass eine Gradmessung aus zwei sehr verschiedenartigen Teilen besteht, einem trigonometrischen und einem astronomischen. Ohne hier weiter ins Einzelne zu gehen, will ich nur bemerken, dass ich den trigonometrischen Teil (der 25 Hauptdreiecke von Göttingen bis Hamburg umfasst) in den Jahren 1821, 1822, 1823 und den astronomischen im Jahr 1827

\par\clearpage
\noindent\null\hfill{\small\textbf{ZUR HANNOVERSCHEN TRIANGULATION.}}\hfill\makebox[0pt][r]{\fontsize{13}{16}\selectfont 419}\par\medskip

\noindent ganz abgeschlossen habe und dass der letztere sowie die Endergebnisse in einem 1828 von mir herausgegebenen Werk »Bestimmung des Breitenunterschiedes zwischen den Sternwarten von Göttingen und Altona durch Beobachtungen am \textsc{Ramsden}schen Zenithsector« bekannt gemacht worden sind. Die Sternwarten von Göttingen und Altona, die durch ein einzigartiges Spiel des Zufalls genau auf demselben Meridian liegen, sind nämlich die Endpunkte des gemessenen Meridianbogens, dessen gesamte Krümmung etwas mehr als zwei Grad beträgt.

\begin{center}
II.\quad Trigonometrische Verbindung der Gradmessungsdreiecke mit den\\
Dreiecken der königlich niederländischen Vermessungen.
\end{center}

\par\noindent\hspace*{\parindent}Der unmittelbare wissenschaftliche Zweck dieser mir durch Erlass vom 8. März 1824 aufgetragenen Erweiterung der trigonometrischen Arbeiten war es, die hannoversche Gradmessung mit der französischen und der englischen zu verbinden. Diese beiden sind bekanntlich miteinander verbunden, und an die französischen Dreiecke schließen sich die mit großer Sorgfalt gemessenen Dreiecke in den Niederlanden an. In der Zeit, als Ostfriesland mit Holland vereinigt war, waren sie bis an die östliche Grenze dieses Fürstentums, weiter südlich dagegen bis nach Bentheim fortgeführt worden. Ich hatte also die Wahl zwischen zwei Wegen, auf denen sich der angestrebte Zweck erreichen ließ, nämlich entweder von den nördlichsten Dreiecken der Gradmessung über Bremen nach Jever oder von den mittleren durch Westfalen nach Bentheim. Der erste Weg wurde unter anderem deshalb vorgezogen, weil sich dadurch zugleich eine Verbindung mit der Nordsee und damit die Bestimmung der absoluten Höhen aller Dreieckspunkte über dem Meeresspiegel erreichen ließ. Diese Dreiecksmessung habe ich in den Jahren 1824 und 1825 ausgeführt und damit die Zahl aller Hauptdreiecke auf 38[\textsuperscript{*})] gebracht.

\par\medskip
\par\noindent\hspace*{\parindent}Zu diesen trigonometrischen Messungen von 1821---1825 habe ich noch einige Bemerkungen hinzuzufügen, da sie sich in mehrerer Hinsicht von ähnlichen Arbeiten unterscheiden, die früher in anderen Ländern ausgeführt wurden.

\par\smallskip\noindent\rule{.24\linewidth}{.4pt}\par\smallskip
{\footnotesize\noindent\hspace*{\parindent}[\textsuperscript{*}) Die Anzahl der unabhängigen Dreiecke der Gradmessung zwischen Göttingen und Hamburg beträgt 21, die Anzahl aller unabhängigen Dreiecke der Gradmessung und ihrer Fortsetzung nach Jever (ohne das Dreieck Hohehagen-Brocken-Inselsberg) 42. Vgl. S. 297.]\par}

\par\smallskip\noindent\hfill{\footnotesize 53\textsuperscript{*}}\hspace*{.12\linewidth}\par

\par\clearpage
\noindent\makebox[0pt][l]{\fontsize{13}{16}\selectfont 420}\hfill{\small\textbf{NACHLASS.}}\hfill\null\par\medskip

\par\noindent\hspace*{\parindent}1) Durch die Anwendung der von mir zuerst eingeführten \textls[120]{Heliotrope} wurden besondere Signaltürme entbehrlich, Dreiecke von einer früher nicht praktikablen Größe möglich und eine ohne dieses Hilfsmittel unerreichbare Genauigkeit der Messungen erreichbar.

\par\noindent\hspace*{\parindent}2) Bei der Ausführung der Messungen habe ich mich nicht auf das für den unmittelbaren Zweck Erforderliche beschränkt (das, wie bereits erwähnt, zunächst nur einem wissenschaftlichen Ziel diente). Ich habe zugleich versucht, die Messungen für die Geografie des Landes so nutzbar wie möglich zu machen, ohne den unmittelbaren Zweck zu beeinträchtigen. Daher sind auch die meisten Türme im Bereich der Dreieckspunkte sehr genau festgelegt, und zwar sowohl die der Städte im Verlauf der Dreieckskette, wie Göttingen, Hildesheim, Wolfenbüttel, Braunschweig, Hannover, Celle, Lüneburg, Harburg, Hamburg, Buxtehude, Stade, Verden, Bremen, Oldenburg, Varel, Jever, als auch die vieler hundert kleinerer Ortschaften. Das zeigen die Übersichtskarten, die meinen Berichten über die Arbeiten des jeweiligen Jahres beigefügt sind. Damit sind feste Anhaltspunkte und Grundlagen für alle später in den betreffenden Landstrichen vorzunehmenden Detailaufnahmen gewonnen, und diese Bestimmungen erhalten dadurch einen bleibenden Wert.

\par\noindent\hspace*{\parindent}3) Schließlich habe ich diese Ergebnisse mit einem eigenen Verfahren in eine Form gebracht, die für die eben genannte weitere Nutzung wesentliche Vorteile bietet.

\begin{center}
III.\quad Trigonometrische Vermessung der Landesteile,\\
die von den Messungen 1821---1825 nicht erfasst wurden.
\end{center}

\par\noindent\hspace*{\parindent}Die in den Jahren 1821---1825 ausgeführten Arbeiten umfassten, obwohl sie zunächst nicht für diesen Zweck bestimmt waren, eine wirkliche trigonometrische Vermessung eines sehr beträchtlichen Teils des Königreichs Hannover. Die allgemein anerkannten Vorteile einer genauen trigonometrischen Landesvermessung ließen es wünschenswert erscheinen, darüber hinauszugehen. Durch Erlass vom 28. April 1828[\textsuperscript{*})] erhielt ich den Auftrag, die Ausdehnung der trigonometrischen Vermessung auf alle von den früheren Arbeiten noch nicht erfassten Landesteile zu leiten.

\par\smallskip\noindent\rule{.24\linewidth}{.4pt}\par\smallskip
{\footnotesize\noindent\hspace*{\parindent}[\textsuperscript{*}) Der im Gauss-Archiv befindliche Erlass ist auf den 14. April datiert.]\par}

\par\clearpage
\noindent\null\hfill{\small\textbf{ZUR HANNOVERSCHEN TRIANGULATION.}}\hfill\makebox[0pt][r]{\fontsize{13}{16}\selectfont 421}\par\medskip

\par\noindent\hspace*{\parindent}Gleichzeitig wurde eine Detailaufnahme angeordnet. Sie sollte alle Landesteile umfassen, die in der in den achtziger Jahren des vorigen Jahrhunderts ausgeführten Messtischaufnahme des ehemaligen Kurfürstentums Hannover noch nicht enthalten waren. Mit der unmittelbaren Leitung dieser Detailaufnahme wurde Oberst \textsc{Prott} beauftragt. Diese Arbeit steht insofern in direkter Verbindung mit der trigonometrischen Vermessung, als deren Ergebnisse die Grundlage und feste Anhaltspunkte für die Messtischarbeiten liefern. Wie sehr die Messtischarbeiten auf diese Weise an Leichtigkeit und Sicherheit gewinnen, hatte sich bereits ein Jahr zuvor bei einigen vorläufigen Detailaufnahmen im Hildesheimer Gebiet gezeigt, obwohl die Arbeiten damals von Offizieren ausgeführt wurden, denen diese Methode zuvor fremd gewesen war.

\par\noindent\hspace*{\parindent}Was ich hier über die unter meine Leitung gestellte trigonometrische Vermessung zu sagen habe, betrifft das beteiligte Personal, den Umfang der bisher abgeschlossenen und schließlich die noch ausstehenden Messungen.

\begin{center}
A.\quad Das Personal.
\end{center}

\par\noindent\hspace*{\parindent}Für eine trigonometrische Vermessung sind zwei ganz verschiedenartige Arbeiten erforderlich, nämlich die Ausführung der Messungen an den betreffenden Orten im Feld und ihre Verarbeitung zu Ergebnissen durch Kombination und Berechnung am Schreibtisch. Den zweiten Teil der Arbeit habe ich bisher ganz selbst übernommen. Den ersten habe ich den Artillerieoffizieren übertragen, die mir in den Jahren 1821---1825 als Gehilfen zur Seite standen und dabei Gelegenheit hatten, sich sowohl mit der Handhabung der Instrumente als auch mit dem Geist und den Besonderheiten meines Verfahrens vertraut zu machen. Diese Offiziere waren Hauptmann \textsc{Müller}, der damalige Premierleutnant \textsc{Hartmann} und mein ältester Sohn, gegenwärtig Premierleutnant im Artillerieregiment. Leutnant \textsc{Hartmann}, der 1831 den Militärdienst verließ, um eine Stelle an der höheren Gewerbeschule in Hannover anzutreten, und den Titularrang eines Hauptmanns hatte, starb im Jahr 1834.

\par\clearpage
\noindent\makebox[0pt][l]{\fontsize{13}{16}\selectfont 422}\hfill{\small\textbf{NACHLASS.}}\hfill\null\par\medskip

\begin{center}
B.\quad Bisher abgeschlossene Hauptdreiecksmessungen.
\end{center}

\par\noindent\hspace*{\parindent}Hätten sich die genannten drei Offiziere jedes Jahr während der ganzen geeigneten Jahreszeit ununterbrochen und ausschließlich dieser Arbeit widmen können, so wäre sie längst abgeschlossen. Verschiedene Hindernisse, die sich teils aus den Dienstverhältnissen der Offiziere, teils aus anderen äußeren Umständen ergaben, führten aber dazu, dass sie in manchen Jahren nur kürzere Zeit, in anderen gar nicht daran arbeiten konnten. Zu dem bereits erwähnten Tod des Hauptmanns \textsc{Hartmann} kam noch hinzu, dass die fortschreitende Detailaufnahme weiterhin besondere Vorbereitungsmessungen erforderte, vor allem in Gegenden, in denen die Kirchtürme weiter verstreut liegen. Dafür konnte jeweils nur einer der genannten Offiziere eingesetzt werden, dessen Zeit dann der Hauptarbeit entzogen wurde. Es folgt eine kurze Übersicht über das bisher in den einzelnen Jahren Geleistete.

{\hangindent=.075\linewidth\hangafter=1\noindent\makebox[.075\linewidth][l]{1828}wurde die Arbeit im Spätsommer begonnen. Hauptmann \textsc{Müller} und Leutnant \textsc{Gauss} schlossen die trigonometrische Vermessung des Eichsfelds größtenteils ab, während Leutnant \textsc{Hartmann} im Amt Hunnesrück und in einem Teil des Hildesheimer Gebiets Vorbereitungsmessungen für die Detailaufnahme vornahm.\par}

{\hangindent=.075\linewidth\hangafter=1\noindent\makebox[.075\linewidth][l]{1829.}\enspace Leutnant \textsc{Gauss} nahm zunächst einige Ergänzungsmessungen im Eichsfeld vor. Den größten Teil des Jahres widmete er jedoch der Triangulation in Westfalen, wo er von der Weser bis Bentheim ein Dreieckssystem ausführte. Hauptmann \textsc{Müller} hatte sich nur kurz an dieser Arbeit beteiligen können, da er bald nach Beginn wegen anderer Aufträge von der Vermessung abberufen wurde. Leutnant \textsc{Hartmann} nahm in diesem Jahr Vorbereitungsmessungen für die Detailaufnahme vor, zuerst im Hildesheimer Gebiet, danach in den Ämtern Uchte, Freudenberg und Auburg.\par}

{\hangindent=.075\linewidth\hangafter=1\noindent\makebox[.075\linewidth][l]{1830.}\enspace Trigonometrische Vermessung des östlichen Teils des Lüneburger Gebiets durch Hauptmann \textsc{Müller} und Leutnant \textsc{Gauss}. Weitere Fortführung der im Vorjahr durch Westfalen gemessenen Dreieckskette bis Ostfriesland durch Leutnant \textsc{Hartmann}.\par}

{\hangindent=.075\linewidth\hangafter=1\noindent\makebox[.075\linewidth][l]{1831}konnte den Arbeiten nur kurze Zeit gewidmet werden. Leutnant \textsc{Gauss}\par}

\par\clearpage
\noindent\null\hfill{\small\textbf{ZUR HANNOVERSCHEN TRIANGULATION.}}\hfill\makebox[0pt][r]{\fontsize{13}{16}\selectfont 423}\par\medskip

{\leftskip=.075\linewidth\noindent schloss die Messungen im Lüneburger Gebiet ab, Leutnant \textsc{Hartmann} die in Ostfriesland. Hauptmann \textsc{Müller} war gar nicht beteiligt.\par}

{\hangindent=.075\linewidth\hangafter=1\noindent\makebox[.075\linewidth][l]{1832}fielen die Messungen ganz aus.\par}

{\hangindent=.075\linewidth\hangafter=1\noindent\makebox[.075\linewidth][l]{1833.}\enspace Triangulation des Landstrichs längs der Weser von der Gegend um Nienburg bis Holzminden durch Hauptmann \textsc{Müller} und Leutnant \textsc{Gauss}. Triangulation des Harzes durch Hauptmann \textsc{Hartmann}.\par}

{\hangindent=.075\linewidth\hangafter=1\noindent\makebox[.075\linewidth][l]{1834}konnte allein Leutnant \textsc{Gauss} den Arbeiten eine, wenn auch sehr kurze, Zeit widmen, teils für Vorbereitungsmessungen zur Detailaufnahme im Osnabrücker Gebiet, teils zur Ergänzung der Messungen des Vorjahres an der Weser.\par}

{\hangindent=.075\linewidth\hangafter=1\noindent\makebox[.075\linewidth][l]{1835.}\enspace Auch in diesem Jahr konnte sich nur Leutnant \textsc{Gauss} für kurze Zeit weiteren Vorbereitungsmessungen im Osnabrücker Gebiet widmen, diesmal mit Unterstützung von Dr. \textsc{Goldschmidt}, Observator an der Göttinger Sternwarte.\par}

{\leftskip=.075\linewidth\noindent In den beiden folgenden Jahren konnte sich allein Hauptmann \textsc{Müller} den Messungen widmen, und zwar\par}

{\hangindent=.075\linewidth\hangafter=1\noindent\makebox[.075\linewidth][l]{1836}der Triangulation des Landstrichs an der Oberweser zwischen Uslar, Göttingen und Münden,\par}

{\hangindent=.075\linewidth\hangafter=1\noindent\makebox[.075\linewidth][l]{1837}weiteren Vorbereitungen für die Detailaufnahme im Osnabrücker Gebiet.\par}

{\leftskip=.075\linewidth\noindent Außerdem erkundete Hauptmann \textsc{Müller} in diesem Jahr vorläufig die Allergegend für ihre künftige trigonometrische Vermessung.\par}

\begin{center}
C.\quad Noch ausstehende Hauptdreiecksmessungen.
\end{center}

\par\noindent\hspace*{\parindent}Die bisher aufgezählten trigonometrischen Vermessungen hängen alle miteinander zusammen. Um das \textls[120]{ganze} Königreich zu erfassen und ein im Wesentlichen vollständiges Ganzes zu bilden, fehlt nur noch die trigonometrische Vermessung zweier größerer Landesteile, nämlich erstens der Gegend rechts und links der Aller und zweitens des nördlichen Teils des Bremer Gebiets. Im ersten Landesteil wurde, wie bereits erwähnt, eine vorläufige Erkundung durchgeführt.

\par\noindent\hspace*{\parindent}Bei Arbeiten dieser Art ist es unmöglich, den erforderlichen Zeit- und Kostenaufwand einigermaßen genau im Voraus zu bestimmen, da der schnellere oder langsamere Fortschritt von vielerlei Umständen abhängt, die teils im Voraus nicht

\par\clearpage
\noindent\makebox[0pt][l]{\fontsize{13}{16}\selectfont 424}\hfill{\small\textbf{NACHLASS.}}\hfill\null\par\medskip

\noindent genau bekannt, teils zufällig sind (z. B. besondere Schwierigkeiten des Geländes oder die Witterungsverhältnisse). Dennoch ist mit einiger Wahrscheinlichkeit anzunehmen, dass Hauptmann \textsc{Müller} die Messungen in zwei Sommern ganz oder doch größtenteils abschließen könnte, wenn er sich ausschließlich diesen Messungen widmen soll.

\par\medskip
\par\noindent\hspace*{\parindent}Die eigentlichen Endergebnisse der trigonometrischen Messungen selbst sind die \textls[120]{Zahlen}, die die Lage der bestimmten Objekte (größtenteils Kirchtürme) mit höchster Genauigkeit festlegen und daher einen bleibenden Wert behalten. Nach dem vollständigen Abschluss der Arbeit wird man sie ordnen und zur größeren Sicherheit in mehreren Abschriften hinterlegen können, davon etwa eine in der Sternwarte. Ob sie auch im Druck veröffentlicht werden sollen, wird dann von einer Entscheidung höherer Stelle abhängen. Ich möchte jedoch nicht unerwähnt lassen, dass die \textsc{Papen}sche Karte des Königreichs, die den vollkommensten Arbeiten dieser Art in jeder Hinsicht gleichkommt und von der bereits 22 Blätter (also ein Drittel des Ganzen) erschienen sind, ihre Genauigkeit diesen Ergebnissen verdankt. Leutnant \textsc{Gauss}, der die Gradeinteilung und die genaue Eintragung aller Grundpositionen für diese Karte übernommen hat, stehen für diesen Zweck alle bisherigen Ergebnisse zur Verfügung.

\par\medskip
\par\noindent\hspace*{\parindent}Abschließend möchte ich noch etwas über die Vorbereitungsmessungen für die Detailaufnahme hinzufügen. Ihr Abschluss lässt sich deshalb gar nicht im Voraus bestimmen, weil sie der Natur der Sache nach notwendig noch beinahe ebenso lange dauern müssen wie die Detailaufnahme selbst. Denn in Gegenden, in denen die durch die Hauptdreiecksmessungen bestimmten Kirchtürme zu weit auseinanderliegen und nicht zahlreich genug sind, um für jedes Messtischblatt genügend feste Anhaltspunkte zu liefern, müssen besondere Signalpfähle gesetzt und ihre Standorte durch kleinere Dreiecke bestimmt werden. Solche Signalpfähle lassen sich aber nicht so zuverlässig schützen, dass man mit Sicherheit damit rechnen könnte, dass sie viele Jahre bestehen bleiben. Da nun zwar nicht unbedingt durch den Verlust eines einzelnen Signalpfahls hier oder dort, wohl aber

\par\clearpage
\noindent\null\hfill{\small\textbf{ZUR HANNOVERSCHEN TRIANGULATION.}}\hfill\makebox[0pt][r]{\fontsize{13}{16}\selectfont 425}\par\medskip

\noindent durch das Verschwinden mehrerer vor ihrer Benutzung für die Messtischaufnahme die ganze für ihre Bestimmung aufgewandte Arbeit verloren wäre, dürfen solche Vorbereitungsarbeiten immer erst höchstens ein Jahr oder ein paar Jahre vor der Detailaufnahme der betreffenden Gegend durchgeführt werden.

\par\noindent\hspace*{\parindent}Im Übrigen ist die Detailaufnahme selbst im Hildesheimer Gebiet, im Eichsfeld, im Amt Hunnesrück sowie in den Ämtern Freudenberg, Uchte und Auburg abgeschlossen. Im Fürstentum Osnabrück ist sie bereits recht weit fortgeschritten. Einen genaueren und vollständigen Bericht darüber wird aber nur Oberst \textsc{Prott} geben können, in dessen Zuständigkeitsbereich diese Arbeit fällt.

\par\noindent\hspace*{\parindent}Göttingen, 8. Februar 1838.

\par\medskip
\begin{center}
[Aus dem Bericht vom 5. Juli 1840 ȟber die trigonometrischen\\
Vermessungen« im Jahr 1839.]
\end{center}

\par\noindent\hspace*{\parindent}\mbox{.\ .\ .\ .\ .\ .}\quad Bei der Beurteilung der Dreieckssysteme darf man nicht übersehen, dass ursprünglich keine allgemeine Landesvermessung beabsichtigt war, sondern zunächst nur eine Gradmessung von Göttingen bis Holstein und danach zunächst eine Erweiterung des Dreieckssystems zum Anschluss an die \textsc{Krayenhoff}schen Dreiecke bis Ostfriesland. Entsprechend \textls[120]{diesen} Zwecken wurden die von mir selbst 1821---1825 gemessenen Dreiecke vom Inselsberg bis Jever \mbox{ausgewählt.\enspace\mbox{.\ .\ .\ .}}\quad Die übrigen, die später hinzukamen, erscheinen als Abzweigungen jener Hauptdreiecke. Infolge dieser Entstehungsweise erscheint das Ganze nicht überall in dem Maße als ein abgerundetes Ganzes aus einem Guss, wie es der Fall gewesen wäre, wenn man dies von Anfang an hätte berücksichtigen müssen. Der eigentliche Zweck, nämlich die genaue Festlegung der wichtigsten dafür geeigneten Punkte im ganzen Land, ist deshalb aber nicht weniger gut \mbox{erreicht.\enspace\mbox{.\ .\ .\ .\ .\ .}}

\par\medskip
\begin{center}
[Aus einem Bericht, Dezember 1844, »über die im Jahre 1844 ausgeführten\\
trigonometrischen Messungen.«]
\end{center}

\par\noindent\hspace*{\parindent}\mbox{.\ .\ .\ .\ .\ .}\quad Ich erlaube mir, noch einige Worte zu allen Messungen der vergangenen Jahre hinzuzufügen.

\par\smallskip\noindent\hspace*{.03\linewidth}{\footnotesize IX.}\hfill{\footnotesize 54}\hspace*{.12\linewidth}\par

\par\clearpage
\noindent\makebox[0pt][l]{\fontsize{13}{16}\selectfont 426}\hfill{\small\textbf{NACHLASS.}}\hfill\null\par\medskip

\par\noindent\hspace*{\parindent}\mbox{.\ .\ .\ .\ .\ .}\quad Die Ergebnisse [d. h. die Koordinaten] wurden nach der Verarbeitung der Messungen jedes Jahr in Verzeichnisse eingetragen. Es gibt sechzehn solcher Teilverzeichnisse, die zusammen etwas mehr als 3000 Bestimmungen enthalten. Die Zahl der Punkte selbst dürfte jedoch um etwa ein Siebtel kleiner sein, da viele Punkte, die in einem späteren Jahr nach dem Hinzukommen neuer Daten genauer oder zuverlässiger bestimmt werden konnten, in mehr als einem Verzeichnis vorkommen. Im ganzen Königreich dürften nur wenige Kirchtürme unbestimmt geblieben sein.

\par\noindent\hspace*{\parindent}Dass diese Verzeichnisse die Grundlage aller seit 16 oder 17 Jahren vorgenommenen Detailaufnahmen und der \textsc{Papen}schen Karten waren, braucht hier nicht weiter ausgeführt zu werden. Von großer Bedeutung ist aber, dass diese Zahlen, die ihren Wert behalten, solange die Gegenstände bestehen, nicht verloren gehen können. Die erwähnten Verzeichnisse werden in der Sternwarte aufbewahrt. Abschriften besitzt auch Leutnant \textsc{Gauss}, der alle Stammpunkte in die \textsc{Papen}schen Karten eingetragen hat. Zur größeren Sicherheit und zur leichteren Benutzung habe ich jetzt begonnen, die Teilverzeichnisse zu einem einzigen zusammenzufassen, das demnach etwa 2600 Punkte enthalten wird.

\par\noindent\hspace*{\parindent}Später könnte es vielleicht für ratsam gehalten werden, dieses Verzeichnis oder einen Auszug daraus im Druck zu veröffentlichen. Für den Augenblick würde ich das aber aus mehreren Gründen noch für verfrüht halten.

\par\noindent\hspace*{\parindent}Erstens, weil sich eine \textls[120]{wissenschaftlich} ausreichende Darstellung der Bedeutung und der vielseitigen Verwendung dieser Zahlen nur nach und nach zusammen mit der Entwicklung meiner eigenen mathematischen Theorien geben lässt, die ich in einer Reihe einzelner Abhandlungen (etwa drei oder vier) vorlegen möchte. Die erste davon ist bereits als Teil des demnächst erscheinenden Bandes der Denkschriften der hiesigen Sozietät der Wissenschaften abgedruckt[\textsuperscript{*})] und auch einzeln in den Buchhandel gebracht worden. Die anderen werde ich nach und nach so bald wie möglich folgen lassen.

\par\noindent\hspace*{\parindent}Zweitens, weil die Zahlen des Verzeichnisses, obwohl sie für jede praktische Verwendung hinreichend, ja mehr als hinreichend genau sind, doch für die den strengsten

\par\smallskip\noindent\rule{.24\linewidth}{.4pt}\par\smallskip
{\footnotesize\noindent\hspace*{\parindent}[\textsuperscript{*}) Untersuchungen über Gegenstände der höhern Geodäsie. Erste Abhandlung. Der Königl. Societät überreicht am 23. October 1843. Band IV, S. 259/300.]\par}

\par\clearpage
\noindent\null\hfill{\small\textbf{ZUR HANNOVERSCHEN TRIANGULATION.}}\hfill\makebox[0pt][r]{\fontsize{13}{16}\selectfont 427}\par\medskip

\noindent theoretischen Anforderungen entsprechende Zusammenführung der verschiedenen Messungen zu einem System noch eine gewisse, wenn auch für sich genommen sehr geringe, Verfeinerung und Nachbesserung zulassen.

\par\noindent\hspace*{\parindent}Drittens, weil ich die Lage der wichtigsten Punkte, insbesondere der Kirchtürme in Städten, neben der Angabe in Koordinaten gern auch in einer anderen Form angeben möchte, nämlich nach der geografischen Breite und Länge. Diese stets mit einem beträchtlichen Zeitaufwand verbundene Umrechnung wird sich erst nach und nach durchführen \mbox{lassen.\quad\mbox{.\ .\ .\ .\ .\ .}}

\par\medskip
\begin{center}
[3.]\\
{\large Hauptdreieckspunkte der hannoverschen Messungen.}
\end{center}

\par\smallskip
{\centering
\setlength{\tabcolsep}{0pt}
\newcommand{\GaussIXCoordinate}[7]{%
  \makebox[\linewidth][c]{%
    \makebox[1.8em][r]{\(#1\)\makebox[0pt][l]{\(\vphantom{0}#5\)}}%
    \hspace{.2em}\makebox[1.35em][r]{\(#2\)\makebox[0pt][l]{\(\vphantom{0}#6\)}}%
    \hspace{.2em}\makebox[1.55em][r]{\(#3\)\makebox[0pt][l]{\(\vphantom{0}#7\)}}%
    \( {,}#4 \)%
  }%
}
\begin{tabular}{@{\vrule width .8pt}p{.34\linewidth}|p{.23\linewidth}|p{.27\linewidth}@{\vrule width .8pt}}
\noalign{\hrule height .8pt}
\rule[-1.0\baselineskip]{0pt}{2.6\baselineskip}\makebox[\linewidth][c]{Dreieckspunkt} &
\makebox[\linewidth][c]{Breite} &
\makebox[\linewidth][c]{Länge v.\ Gött.\ Merid.} \\
\hline
\multicolumn{1}{@{\vrule width .8pt}c|}{\rule{0pt}{.45\baselineskip}} &
& \\
\hspace*{.6em}Inselsberg, hess. Dr.-P. & \GaussIXCoordinate{50}{51}{8}{619}{^{0}}{^{\prime}}{^{\prime\prime}} & \GaussIXCoordinate{-0}{31}{24}{311}{^{0}}{^{\prime}}{^{\prime\prime}} \\
\hspace*{.6em}Hohehagen & \GaussIXCoordinate{51}{28}{31}{234}{}{}{} & \GaussIXCoordinate{+0}{10}{45}{143}{}{}{} \\
\hspace*{.6em}Sternwarte in Göttingen & \GaussIXCoordinate{51}{31}{48}{028}{}{}{} & \makebox[\linewidth][c]{\(0\)} \\
\hspace*{.6em}Meridianzeichen & \GaussIXCoordinate{51}{34}{30}{309}{}{}{} & \GaussIXCoordinate{-0}{0}{0}{007}{}{}{} \\
\hspace*{.6em}Brocken & \GaussIXCoordinate{51}{48}{1}{849}{}{}{} & \GaussIXCoordinate{-0}{40}{23}{071}{}{}{} \\
\hspace*{.6em}Hils & \GaussIXCoordinate{51}{53}{53}{000}{}{}{} & \GaussIXCoordinate{+0}{6}{41}{152}{}{}{} \\
\hspace*{.6em}Lichtenberg & \GaussIXCoordinate{52}{7}{22}{020}{}{}{} & \GaussIXCoordinate{-0}{20}{33}{348}{}{}{} \\
\hspace*{.6em}Deister & \GaussIXCoordinate{52}{14}{5}{778}{}{}{} & \GaussIXCoordinate{+0}{20}{35}{702}{}{}{} \\
\hspace*{.6em}Garssen & \GaussIXCoordinate{52}{39}{40}{397}{}{}{} & \GaussIXCoordinate{-0}{12}{19}{154}{}{}{} \\
\hspace*{.6em}Scharnhorst & \GaussIXCoordinate{52}{43}{43}{000}{}{}{} & \GaussIXCoordinate{-0}{19}{1}{620}{}{}{} \\
\hspace*{.6em}Breithorn & \GaussIXCoordinate{52}{49}{45}{985}{}{}{} & \GaussIXCoordinate{-0}{19}{11}{230}{}{}{} \\
\hspace*{.6em}Falkenberg & \GaussIXCoordinate{52}{50}{52}{481}{}{}{} & \GaussIXCoordinate{+0}{4}{34}{858}{}{}{} \\
\hspace*{.6em}Hauselberg & \GaussIXCoordinate{52}{51}{36}{513}{}{}{} & \GaussIXCoordinate{-0}{14}{28}{226}{}{}{} \\
\hspace*{.6em}Elmhorst & \GaussIXCoordinate{52}{59}{42}{886}{}{}{} & \GaussIXCoordinate{+0}{18}{24}{543}{}{}{} \\
\hspace*{.6em}Steinberg & \GaussIXCoordinate{52}{59}{54}{987}{}{}{} & \GaussIXCoordinate{+0}{40}{7}{347}{}{}{} \\
\hspace*{.6em}Bottel & \GaussIXCoordinate{53}{2}{22}{919}{}{}{} & \GaussIXCoordinate{+0}{39}{22}{916}{}{}{} \\
\hspace*{.6em}Wulfsode & \GaussIXCoordinate{53}{4}{15}{113}{}{}{} & \GaussIXCoordinate{-0}{17}{48}{837}{}{}{} \\
\hspace*{.6em}Bremen, Ansgar.-Turm & \GaussIXCoordinate{53}{4}{48}{867}{}{}{} & \GaussIXCoordinate{-1}{8}{22}{781}{}{}{} \\
\multicolumn{1}{@{\vrule width .8pt}c|}{\rule{0pt}{.45\baselineskip}} &
& \\
\end{tabular}\par
}

\par\smallskip\noindent\hfill{\footnotesize 54\textsuperscript{*}}\hspace*{.12\linewidth}\par
